\chapter{Những Điều Về Tính Cách}
\markboth{Những Điều Về Tính Cách}{Things a Teacher Wants to Say About Character}

\section{Trung Thực Trong Học Tập}
\begin{truyen}{Trung Thực Trong Học Tập}{Be Honest in Learning}
\chuhoa{T}{rung thực trong học tập không chỉ là không quay cóp. Nó còn là không giả vờ hiểu khi mình chưa hiểu, không giả vờ cố gắng khi mình đang bỏ bê.}
\textit{(Honesty in learning is not only about not cheating. It is also about not pretending to understand, and not pretending to try when you are neglecting the work.)}

Có thể em vẫn qua được một bài kiểm tra bằng sự không trung thực.
\textit{(You may still pass a test through dishonesty.)}

Nhưng em sẽ khó xây được một con người vững bằng cách đó.
\textit{(But it is hard to build a steady person that way.)}

Trung thực làm em có thể đau một lúc, nhưng nó giữ cho em còn tôn trọng được chính mình.
\textit{(Honesty may hurt for a moment, but it helps you keep respect for yourself.)}
\end{truyen}

\begin{baihoc}
Học cho thật luôn chậm hơn ăn gian một chút, nhưng nó cho em một nền vững hơn rất nhiều.
\textit{(Learning honestly is slower than cheating, but it gives you a much stronger foundation.)}
\end{baihoc}

\begin{ghinhoanh}
\textbf{Short English to remember:}

Honest learning builds a stronger self.

\end{ghinhoanh}

\begin{tuhocnhanh}
1. Trung thực trong học tập gồm những gì ngoài chuyện không quay cóp? \textit{What does academic honesty include beyond not cheating?}

2. Vì sao gian dối có thể qua mắt người khác nhưng làm yếu mình? \textit{Why can dishonesty fool others but weaken you?}

3. Em cần trung thực hơn ở chỗ nào trong việc học của mình? \textit{Where do you need more honesty in your learning?}
\end{tuhocnhanh}
\ngancach

\section{Tôn Trọng Thầy Cô}
\begin{truyen}{Tôn Trọng Thầy Cô}{Respect Teachers}
\chuhoa{T}{ôn trọng thầy cô không phải vì người lớn luôn đúng. Nó là vì việc học cần một thái độ biết lắng nghe và biết cư xử cho đàng hoàng.}
\textit{(Respecting teachers is not because adults are always right. It is because learning needs an attitude of listening and proper conduct.)}

Em có thể không đồng ý với một điều gì đó.
\textit{(You may disagree with something.)}

Nhưng cách em không đồng ý cũng nói lên con người em.
\textit{(But the way you disagree also reveals your character.)}

Tôn trọng không làm em thấp đi. Nó làm cách ứng xử của em trưởng thành hơn.
\textit{(Respect does not lower you. It makes your way of behaving more mature.)}
\end{truyen}

\begin{baihoc}
Tôn trọng người dạy không chỉ là phép lịch sự. Nó còn là một phần của văn hóa học tập tử tế.
\textit{(Respecting teachers is not only politeness. It is part of a decent learning culture.)}
\end{baihoc}

\begin{ghinhoanh}
\textbf{Short English to remember:}

Respect makes disagreement more mature.

\end{ghinhoanh}

\begin{tuhocnhanh}
1. Vì sao tôn trọng không đồng nghĩa với đồng ý hết? \textit{Why is respect not the same as complete agreement?}

2. Cách không đồng ý nói gì về tính cách? \textit{What does disagreement style reveal about character?}

3. Em có thể cư xử chín hơn với thầy cô ở điểm nào? \textit{Where can you act more maturely with teachers?}
\end{tuhocnhanh}
\ngancach

\section{Giúp Đỡ Bạn Bè}
\begin{truyen}{Giúp Đỡ Bạn Bè}{Help Your Friends}
\chuhoa{G}{iúp đỡ bạn bè trong học tập không làm kiến thức của em mất đi. Nhiều khi nó còn làm em hiểu sâu hơn.}
\textit{(Helping friends in study does not take your knowledge away. It often helps you understand more deeply.)}

Khi giải thích cho ai đó, em buộc phải sắp xếp lại điều mình biết cho rõ ràng.
\textit{(When explaining to someone, you must organize what you know more clearly.)}

Giúp bạn cũng làm lớp học bớt lạnh hơn.
\textit{(Helping friends also makes the classroom less cold.)}

Tất nhiên, giúp khác với làm hộ. Điều tốt là giúp người khác đứng lên, không phải làm thay để họ khỏi phải cố.
\textit{(Of course, helping is different from doing the work for them. Good help lifts others up rather than replacing their effort.)}
\end{truyen}

\begin{baihoc}
Giúp đỡ đúng cách làm người khác đỡ yếu đi, nhưng cũng làm mình đàng hoàng hơn.
\textit{(Helping in the right way makes others less helpless and makes you more decent as well.)}
\end{baihoc}

\begin{ghinhoanh}
\textbf{Short English to remember:}

Good help lifts people; it does not replace their effort.

\end{ghinhoanh}

\begin{tuhocnhanh}
1. Vì sao giúp bạn cũng làm mình hiểu bài hơn? \textit{Why can helping friends deepen your understanding?}

2. Giúp đỡ và làm hộ khác nhau thế nào? \textit{How is helping different from doing it for someone?}

3. Em có thể giúp một bạn theo cách nào cho đúng? \textit{How can you help a friend in the right way?}
\end{tuhocnhanh}
\ngancach

\section{Biết Xin Lỗi}
\begin{truyen}{Biết Xin Lỗi}{Know How to Apologize}
\chuhoa{X}{in lỗi không làm em yếu đi. Nó cho thấy em đủ trưởng thành để nhận phần sai của mình.}
\textit{(Apologizing does not weaken you. It shows you are mature enough to own your part of the wrong.)}

Nhiều người ngại xin lỗi vì sợ mất mặt.
\textit{(Many people avoid apologizing because they fear losing face.)}

Nhưng thật ra, người không bao giờ nhận lỗi mới dễ làm mình nhỏ đi trong mắt người khác.
\textit{(In truth, those who never admit fault often become smaller in others' eyes.)}

Một lời xin lỗi chân thành không xóa hết hậu quả, nhưng nó mở ra con đường sửa lại điều đã hỏng.
\textit{(A sincere apology does not erase all consequences, but it opens a path toward repair.)}
\end{truyen}

\begin{baihoc}
Xin lỗi là một biểu hiện của tự trọng và trách nhiệm, không phải của yếu đuối.
\textit{(Apologizing is a sign of self-respect and responsibility, not weakness.)}
\end{baihoc}

\begin{ghinhoanh}
\textbf{Short English to remember:}

A sincere apology is a form of courage.

\end{ghinhoanh}

\begin{tuhocnhanh}
1. Vì sao nhiều người ngại xin lỗi? \textit{Why do many people hesitate to apologize?}

2. Một lời xin lỗi tốt cần điều gì? \textit{What makes an apology sincere?}

3. Em có đang nợ ai một lời xin lỗi không? \textit{Do you owe someone an apology right now?}
\end{tuhocnhanh}
\ngancach

\section{Biết Cảm Ơn}
\begin{truyen}{Biết Cảm Ơn}{Know How to Be Grateful}
\chuhoa{C}{ảm ơn là một việc nhỏ nhưng làm lòng người khác ấm lên rất nhiều.}
\textit{(Saying thank you is small, but it can warm another person's heart greatly.)}

Người biết cảm ơn thường nhìn thấy điều tốt mình đang nhận được rõ hơn.
\textit{(People who know gratitude often notice the good they are receiving more clearly.)}

Trong học đường, một lời cảm ơn bạn đã chỉ bài, cảm ơn thầy cô đã kiên nhẫn, cảm ơn gia đình đã lo cho mình đi học là điều không nhỏ.
\textit{(In school life, thanking a friend for explaining, a teacher for patience, or family for supporting school is not a small thing.)}

Cảm ơn không phải xã giao rỗng. Nó là một cách sống ít vô tâm hơn.
\textit{(Gratitude is not empty social ritual. It is a way of living with less indifference.)}
\end{truyen}

\begin{baihoc}
Biết cảm ơn giúp con người bớt cay cú với đời và sống mềm hơn với nhau.
\textit{(Gratitude helps people become less bitter and softer with one another.)}
\end{baihoc}

\begin{ghinhoanh}
\textbf{Short English to remember:}

Gratitude makes the heart less hard.

\end{ghinhoanh}

\begin{tuhocnhanh}
1. Vì sao lời cảm ơn nhỏ mà không hề nhỏ? \textit{Why is a small thank you not actually small?}

2. Người biết cảm ơn thường khác ở điểm nào? \textit{How are grateful people different?}

3. Hôm nay em cần cảm ơn ai? \textit{Whom do you need to thank today?}
\end{tuhocnhanh}
\ngancach

\section{Giữ Lời Hứa}
\begin{truyen}{Giữ Lời Hứa}{Keep Your Word}
\chuhoa{M}{ột người giữ lời chưa chắc nói hay, nhưng rất đáng tin.}
\textit{(A person who keeps their word may not speak beautifully, but they are trustworthy.)}

Trong học tập, giữ lời hứa có thể đơn giản như hứa làm xong bài thì làm xong, hứa giúp bạn thì nhớ giúp, hứa sửa lỗi thì thật sự sửa.
\textit{(In learning, keeping your word may be as simple as finishing promised work, helping a friend as said, or actually correcting a mistake you said you would fix.)}

Người thất hứa mãi sẽ làm người khác mỏi.
\textit{(A person who constantly breaks promises makes others weary.)}

Giữ lời làm em lớn lên trong sự đáng tin, và đó là một tài sản rất lớn trong đời.
\textit{(Keeping your word helps you grow into someone dependable, and that is a major asset in life.)}
\end{truyen}

\begin{baihoc}
Giữ lời là một dạng kỷ luật với chính mình và là món quà đáng tin với người khác.
\textit{(Keeping your word is a discipline toward yourself and a gift of trust toward others.)}
\end{baihoc}

\begin{ghinhoanh}
\textbf{Short English to remember:}

Trust grows where promises are kept.

\end{ghinhoanh}

\begin{tuhocnhanh}
1. Vì sao giữ lời làm con người đáng tin hơn? \textit{Why does keeping your word build trust?}

2. Em hay thất hứa với ai nhất, kể cả với chính mình? \textit{With whom do you break promises most, including yourself?}

3. Một lời hứa nào em cần làm thật hơn từ hôm nay? \textit{What promise do you need to keep more honestly from today?}
\end{tuhocnhanh}
\ngancach

\section{Không Cười Khi Bạn Sai}
\begin{truyen}{Không Cười Khi Bạn Sai}{Do Not Laugh When Someone Is Wrong}
\chuhoa{M}{ột tiếng cười chế nhạo có thể làm bạn mình im rất lâu.}
\textit{(One mocking laugh can keep a friend silent for a long time.)}

Lúc một người đang sai, họ đã đủ ngượng rồi.
\textit{(When someone is already wrong, they are often ashamed enough.)}

Nếu ta còn cười để làm họ nhỏ thêm, ta đang góp phần biến lớp học thành chỗ người ta sợ hơn là chỗ người ta học.
\textit{(If we laugh to make them smaller, we help turn the classroom into a place of fear rather than learning.)}

Người tử tế biết dừng lại trước cái sai của người khác bằng sự chừng mực.
\textit{(A decent person knows how to meet another's mistake with restraint.)}
\end{truyen}

\begin{baihoc}
Đừng làm nỗi xấu hổ của người khác lớn thêm chỉ để mình thấy vui một lúc.
\textit{(Do not enlarge someone else's shame just for a brief moment of amusement.)}
\end{baihoc}

\begin{ghinhoanh}
\textbf{Short English to remember:}

Mocking someone else's mistake can silence their courage.

\end{ghinhoanh}

\begin{tuhocnhanh}
1. Vì sao một tiếng cười có thể làm hại bạn lâu? \textit{Why can one laugh hurt a friend for a long time?}

2. Lớp học sẽ ra sao nếu ai sai cũng bị cười? \textit{What happens when every mistake is mocked?}

3. Em có thể phản ứng tử tế hơn thế nào khi bạn sai? \textit{How can you respond more kindly when a friend is wrong?}
\end{tuhocnhanh}
\ngancach

\section{Học Cách Lắng Nghe}
\begin{truyen}{Học Cách Lắng Nghe}{Learn to Listen}
\chuhoa{L}{ắng nghe là một kỹ năng khó hơn người ta tưởng.}
\textit{(Listening is a harder skill than people often think.)}

Nhiều khi mình chỉ đang chờ tới lượt mình nói, chứ chưa thật sự nghe người kia.
\textit{(Often we are only waiting for our turn to speak, not truly listening.)}

Trong lớp học, biết nghe giúp em hiểu bài hơn. Trong tình bạn, biết nghe giúp em hiểu người hơn.
\textit{(In class, listening helps you understand lessons better. In friendship, it helps you understand people better.)}

Một người biết lắng nghe thường bớt vội phán xét hơn.
\textit{(A person who listens well often judges less hastily.)}
\end{truyen}

\begin{baihoc}
Nghe cho hết là một biểu hiện của tôn trọng và cũng là con đường để hiểu sâu hơn.
\textit{(Listening fully is a form of respect and also a path toward deeper understanding.)}
\end{baihoc}

\begin{ghinhoanh}
\textbf{Short English to remember:}

Listening well is a form of respect.

\end{ghinhoanh}

\begin{tuhocnhanh}
1. Vì sao lắng nghe khó hơn mình nghĩ? \textit{Why is listening harder than it seems?}

2. Biết nghe giúp việc học và tình bạn khác đi thế nào? \textit{How does listening change learning and friendship?}

3. Em có đang thật sự nghe hay chỉ chờ tới lượt mình nói? \textit{Are you truly listening or only waiting to speak?}
\end{tuhocnhanh}
\ngancach

\section{Giữ Bình Tĩnh}
\begin{truyen}{Giữ Bình Tĩnh}{Stay Calm}
\chuhoa{B}{ình tĩnh không có nghĩa là em không có cảm xúc. Nó nghĩa là em không để cảm xúc kéo mình đi quá xa trong lúc cần nghĩ rõ.}
\textit{(Calmness does not mean having no emotions. It means not letting emotion drag you too far when clear thinking is needed.)}

Khi nóng giận, xấu hổ hay hoảng, con người rất dễ nói hoặc làm điều mình sẽ tiếc.
\textit{(When angry, ashamed, or panicked, people easily say or do things they later regret.)}

Giữ bình tĩnh giúp em còn chỗ để chọn phản ứng cho tử tế hơn.
\textit{(Staying calm gives you room to choose a more decent response.)}

Đó là kỹ năng phải tập, không phải lúc nào cũng tự có.
\textit{(It is a skill that must be practiced, not something always available naturally.)}
\end{truyen}

\begin{baihoc}
Bình tĩnh giữ cho cái đầu còn làm việc được ngay cả khi lòng đang động.
\textit{(Calmness keeps the mind working even while the heart is unsettled.)}
\end{baihoc}

\begin{ghinhoanh}
\textbf{Short English to remember:}

Calmness creates room for a better choice.

\end{ghinhoanh}

\begin{tuhocnhanh}
1. Giữ bình tĩnh khác với giấu cảm xúc ở đâu? \textit{How is calmness different from hiding emotion?}

2. Khi mất bình tĩnh, em dễ làm điều gì nhất? \textit{What do you do most easily when losing calm?}

3. Cách nào giúp em chậm lại trước khi phản ứng? \textit{What helps you slow down before reacting?}
\end{tuhocnhanh}
\ngancach

\section{Khiêm Tốn}
\begin{truyen}{Khiêm Tốn}{Stay Humble}
\chuhoa{K}{hiêm tốn không phải tự chê mình. Nó là biết mình có cái tốt nhưng không lấy cái tốt đó để đè người khác.}
\textit{(Humility is not self-belittling. It is knowing your strengths without using them to press others down.)}

Một học sinh giỏi mà kiêu dễ làm người khác ngại gần.
\textit{(A strong student who becomes arrogant often drives others away.)}

Một người biết nhiều mà vẫn còn muốn học thêm mới là người thật sự có chiều sâu.
\textit{(A person who knows much yet still wants to keep learning is one with real depth.)}

Khiêm tốn giữ cho em còn mở trước điều mình chưa biết.
\textit{(Humility keeps you open to what you still do not know.)}
\end{truyen}

\begin{baihoc}
Khiêm tốn làm con người vừa dễ học thêm, vừa dễ sống cùng người khác hơn.
\textit{(Humility makes it easier both to keep learning and to live with others.)}
\end{baihoc}

\begin{ghinhoanh}
\textbf{Short English to remember:}

Humility keeps the mind teachable.

\end{ghinhoanh}

\begin{tuhocnhanh}
1. Khiêm tốn khác với tự ti ở đâu? \textit{How is humility different from insecurity?}

2. Vì sao người biết nhiều vẫn cần khiêm tốn? \textit{Why does a knowledgeable person still need humility?}

3. Em có đang dùng điều tốt của mình để làm người khác nhỏ đi không? \textit{Are you using your strengths to make others feel smaller?}
\end{tuhocnhanh}
\ngancach
